\documentclass[11pt]{article}

\usepackage[margin=1in]{geometry}
\usepackage{amsmath,amssymb,amsthm,mathtools}
\usepackage{enumitem}
\usepackage{microtype}
\usepackage[hidelinks]{hyperref}
\usepackage{xurl}

\newtheorem{definition}{Definition}
\newtheorem{lemma}{Lemma}
\newtheorem{remark}{Remark}

\newcommand{\defeq}{\stackrel{\mathrm{def}}{=}}

\newcommand{\TV}{\mathrm{TV}}
\newcommand{\Law}{\mathsf{Law}}

\newcommand{\EE}{\mathbb{E}}

\title{Probabilistic Separation Budgets for Primary Paths\\
\large Separation failure as a receipt-grade $(\varepsilon,\delta)$ adjunct (Series Synthesis)}
\author{}
\date{Unified archive note v0.06 (2026-02-28; archive v130)}

\begin{document}
\maketitle

\begin{abstract}
Many reader-privacy deployments rely on a \emph{separation} primary path: one role observes client-side metadata while a different role observes destination-side metadata, and the claim is interpreted only if these roles do not collude within the threat window.
Synthesis~22 standardizes a digest-bound separation manifest so this assumption is explicit and diffable.

This note adds a minimal quantitative adjunct: if a deployment is willing to publish an upper bound $\rho$ on separation failure (collusion / handle-join) within the declared threat window, then a ``0-bit'' separation intent yields a guarded \emph{additive} bound
\[
\Pr[X\in E\mid K=k]\;\le\;\Pr[X\in E]+\rho
\qquad\text{for all events $E$ and all $k$}.
\]
Equivalently, $\TV(\Law(X\mid K=k),\Law(X))\le \rho$ for all $k$.
This can be carried as a receipt-grade checkpoint $(0,\rho)$ in the series' existing $(\varepsilon,\delta)$ calculus \emph{without changing the public receipt schema}.
\end{abstract}

\paragraph{Scope.}
This is an interface note.
It does \emph{not} provide a universal way to estimate collusion risk.
It only specifies how a published \emph{upper bound} on separation failure can be carried as a digest-bound adjunct and how it propagates to endpoint interpretation.

\paragraph{Imports.}
Separation manifests are owned by Synthesis~22.\cite{series:primarysepmanifest}
Receipt line-item semantics are owned by Synthesis~12 and equivalence/drift semantics by Synthesis~18.\cite{series:receiptitems,series:planstateequiv}
Threat windows are owned by Synthesis~4.\cite{series:threatwindows}
Endpoint meaning is owned by Anonymity~A and the endpoint bridge.\cite{series:oddsinflation,series:endpointbridge}
Horizon composition and checkpoint accounting are centralized in Synthesis~19.\cite{series:receiptacct}
A concrete instantiation of a separation-guard primary path appears in the worked example (Synthesis~17).\cite{series:workedexample}

\paragraph{Exports.}
A minimal ``probabilistic separation'' adjunct: a digest-bound $\rho$ field bound to a separation manifest and threat window, plus an editorial rule for how to interpret it as a $\delta$ term.

\section{The model}
Fix a guarded exposure surface (ENF-ID) and threat window (TW-ID).\cite{series:traceifaces,series:threatwindows}
Under the separation manifest, the deployment intends that the primary-path observation $X$ is \emph{key independent} on that surface (a ``0-bit'' linkability surface).
In practice, the separation can fail (collusion, shared handles, correlated telemetry joins).

\begin{definition}[Separation failure indicator]
Let $F\in\{0,1\}$ be a (possibly hidden) indicator for whether separation is violated during the threat window for the guarded surface.
A deployment publishes an upper bound
\[\Pr[F=1] \le \rho\]
for some $\rho\in[0,1]$.
The bound is interpreted as being conditioned on the same public state scope declared by \texttt{state\_decl\_id} when applicable.\cite{series:planstateequiv}
\end{definition}

\begin{remark}[Why $\rho$ is receipt-grade]
The series already interprets one-vs-mixture privacy via $(\varepsilon,\delta)$ budgets and composes them via checkpoint calculus.\cite{series:oddsinflation}
Publishing $\rho$ therefore does not introduce a new metric family; it introduces a \emph{failure-probability adjunct} that deterministically maps to the checkpoint $(0,\rho)$ under the same endpoint semantics.
\end{remark}

\section{From probabilistic separation to an $(\varepsilon,\delta)$ bound}
The key observation is: if separation holds, the guarded surface is key independent; if it fails, we assume nothing.
This yields an immediate additive bound whose parameters are determined solely by $\rho$.

\begin{lemma}[$\rho$-separation yields a guarded $(0,\rho)$ checkpoint]
\label{lem:rho}
Let $K$ be the hidden key (or destination class) and let $X$ be the declared primary-path observation on the guarded exposure surface.
Let $F\in\{0,1\}$ indicate separation failure during the threat window.
Assume:
(i) conditioned on $F=0$, the law of $X$ is independent of $K$ (the ``0-bit'' separation intent), and
(ii) for all $k$, $\Pr[F=1\mid K=k]\le \rho$.
Then for any prior on $K$ and any event $E$,
\[
\Pr[X\in E\mid K=k] \le \Pr[X\in E] + \rho
\qquad\text{for all }k,
\]
and therefore $\TV(\Law(X\mid K=k),\Law(X))\le \rho$ for all $k$.
\end{lemma}

\begin{proof}
Fix an event $E$ and define
$a\defeq \Pr[X\in E\mid F=0]$,
$\rho_k\defeq \Pr[F=1\mid K=k]$, and
$b_k\defeq \Pr[X\in E\mid K=k,F=1]$.
Key-independence under $F=0$ makes $a$ independent of $k$.
Then
\[
\Pr[X\in E\mid K=k] = (1-\rho_k)a + \rho_k b_k = a + \rho_k(b_k-a).
\]
Let $t_k\defeq \rho_k(b_k-a)$.
Since $b_k\in[0,1]$ and $\rho_k\le\rho$, we have
$t_k\in[-\rho a,\rho(1-a)]$.
Let $\bar t\defeq \EE[t_K]$ under the prior on $K$.
Then $\Pr[X\in E]=a+\bar t$, so
\[
\Pr[X\in E\mid K=k]-\Pr[X\in E] = t_k-\bar t \le \max t - \min t \le \rho.
\]
\end{proof}

\begin{remark}[On multiplicative displays]
Lemma~\ref{lem:rho} is already an $(\varepsilon,\delta)$-style statement with $\varepsilon=0$ and $\delta=\rho$.
If a consumer insists on writing a multiplicative factor, note that $(0,\rho)$ trivially implies $(\varepsilon,\rho)$ for any $\varepsilon\ge 0$.
There is no canonically ``best'' $\varepsilon$ without additional assumptions.
\end{remark}

\begin{remark}[Endpoint meaning]
Lemma~\ref{lem:rho} is a \emph{guarded} statement: with probability at most $\rho$, the separation claim may fail completely.
Downstream endpoint interpretation should therefore treat $(0,\rho)$ as an additive checkpoint in the same style as approximate DP.\cite{series:oddsinflation}
Publishing a larger bound $\rho'\ge\rho$ is a monotone \emph{relaxation} (weaker guarantee); publishing a smaller bound is a \emph{tightening} that should be backed by evidence/replay.
\end{remark}

\section{How to publish $\rho$ without changing the schema}
Synthesis~12 deliberately keeps the receipt schema minimal.\cite{series:receiptitems}
Accordingly, the probabilistic separation adjunct is carried as a digest-bound field adjacent to the receipt:
\begin{itemize}[leftmargin=*]
\item The primary-path line item includes the \texttt{sep\_manifest\_id} (or equivalent manifest digest pointer) as described in Synthesis~22.\cite{series:primarysepmanifest}
\item The same line item carries \texttt{rho\_upper} inside \texttt{obs\_model} (or inside \texttt{knobs}) and tags whether it is a \emph{policy bound} or an \emph{evidence-backed estimate}.
\item Drift rule (aligned to Synthesis~18): treat any \emph{tightening} of \texttt{rho\_upper} as replay/evidence-required; treat a \emph{relaxation} ($\rho'\ge\rho$) as a declared \texttt{lineage=coarsening} change that does not require new evidence but must still be published and diffed.\cite{series:planstateequiv}
\end{itemize}

\paragraph{Interpretation rule.}
Auditors and endpoint interpreters treat \texttt{rho\_upper} as defining the guarded checkpoint $(0,\rho)$ via Lemma~\ref{lem:rho}.
\section{Research warranted (what this note does not solve)}
This note assumes a published upper bound $\rho$.
Estimating $\rho$ in a way that is (i) meaningful, (ii) replayable, and (iii) resistant to adversarial ``audit$\,\neq\,$live'' bias is an open engineering problem.
A plausible publication shape is a conditional certificate (Synthesis~21) whose evidence object is an \emph{evidence bundle} that makes separation failure mechanically checkable.\cite{series:condcert,series:auditorreplay}

\subsection{Toy bounds (when you are willing to assume a model)}
This note is deliberately not a ``collusion-risk estimator,'' but it is still useful to record a few \emph{model-to-$\rho$} translations that are short enough to review.
These are only meaningful if the corresponding assignment/rotation policy is committed in the separation manifest (Synthesis~22).\cite{series:primarysepmanifest}

\begin{lemma}[Independent compromise model $\Rightarrow$ multiplicative $\rho$ bound]
Assume the separation manifest pins a policy in which a lookup chooses a relay instance $R$ from a declared relay set and a router instance $U$ from a declared router set, and the operator sets are declared disjoint.
If the relay-side compromise event $A_R$ and router-side compromise event $A_U$ satisfy
\(\Pr[A_R]\le p_R\), \(\Pr[A_U]\le p_U\), and \(A_R\) is independent of \(A_U\),
then the separation-failure probability (collusion / join of both sides) is bounded by
\[\rho \;\le\; \Pr[A_R\wedge A_U] \;\le\; p_R p_U.
\]
\end{lemma}

\begin{remark}[Finite-adversary clause as a special case]
If the manifest asserts ``at most $t_R$ of $m_R$ relay operators are compromised'' and the relay is selected uniformly at random from that set, then one may take $p_R=t_R/m_R$ (and similarly $p_U=t_U/m_U$) in the lemma above.
This is not a theorem about real adversaries; it is simply the translation from a committed assignment policy and a committed adversary-count clause into an upper bound.
\end{remark}

\begin{remark}[Threat-window aggregation]
If the threat window covers $d$ independent (or union-bounded) epochs with per-epoch separation failure bounds $\rho_1,\dots,\rho_d$, then
\(\Pr[\exists i: F_i=1] \le \sum_i \rho_i\).
In particular, under a uniform per-epoch bound $\rho_\mathrm{epoch}$ one may publish the conservative window bound $\rho\le \min\{1,\,d\,\rho_\mathrm{epoch}\}$.
This is the correct place to encode ``rotation helps'' in a receipt-grade way: the manifest pins the rotation policy, and the published $\rho$ is the resulting window-level bound.
\end{remark}

\begin{itemize}[leftmargin=*]
\item \textbf{Handle-join and key-join detectors:} a logged set of forbidden joins (shared identifiers, shared TLS client certs, shared routing keys) over the threat window, with the detector definition digest-bound in the plan spec.
\item \textbf{Administrative non-collusion attestations:} signed statements binding relay/router operators, billing accounts, and incident-response channels (useful but not sufficient alone).
\item \textbf{Cryptographic separation anchors:} evidence that relay- and router-side keys are rooted in disjoint HSM domains or disjoint issuing CAs for the window, so ``shared handle'' joins are detectable.
\item \textbf{Relay/router log cross-checks:} consistency checks that the sets of requests/epochs seen by each role match the declared interface without revealing per-request linkage (e.g., via digest commitments to per-epoch aggregates).
\item \textbf{Audit$\,\neq\,$live hardening:} replay plans that sample both audit and live traffic to detect conditional failures (e.g., collusion only when audit flags are absent).
\end{itemize}

\begin{thebibliography}{99}\small

\bibitem{series:primarysepmanifest}
Anonymous.
\newblock Primary-Path Separation Manifests: Guard Objects for Non-Collusion Assumptions in Reader-Privacy Receipts.
\newblock Synthesis~22, The-One-True-Archive (2026).

\bibitem{series:receiptitems}
Anonymous.
\newblock Receipt Line Items for Anonymity Claims: A Minimal Tuple Schema for Published Budgets.
\newblock Synthesis~12, The-One-True-Archive (2026).

\bibitem{series:planstateequiv}
Anonymous.
\newblock Digest-Bound Replay Plans and State Declarations: Equivalence, Minimality, and Change Control.
\newblock Synthesis~18, The-One-True-Archive (2026).

\bibitem{series:threatwindows}
Anonymous.
\newblock Threat Models and Repetition Windows for Anonymous DHTs.
\newblock Synthesis~4, The-One-True-Archive (2026).

\bibitem{series:traceifaces}
Anonymous.
\newblock Trace Interfaces for Anonymous {DHT}s.
\newblock Synthesis~3, The-One-True-Archive (2026).

\bibitem{series:oddsinflation}
Anonymous.
\newblock Posterior-Mass Inflation Anonymity: Pointwise Maximal Leakage, Composition, and the True-Odds Boundary.
\newblock Anonymity series Paper~A, The-One-True-Archive (2026).

\bibitem{series:endpointbridge}
Anonymous.
\newblock Endpoint Metrics Bridge for Anonymous Overlays: Posterior-Mass Inflation, PML, MaxL, and True Odds.
\newblock Synthesis~5, The-One-True-Archive (2026).

\bibitem{series:workedexample}
Anonymous.
\newblock Worked Example: Instantiating a Tiered Reader-Privacy Receipt.
\newblock Synthesis~17, The-One-True-Archive (2026).

\bibitem{series:condcert}
Anonymous.
\newblock Conditional Per-Step Budget Certificates: A Receipt-Grade Interface for Adaptive Horizon Claims.
\newblock Synthesis~21, The-One-True-Archive (2026).

\bibitem{series:auditorreplay}
Anonymous.
\newblock Auditor Replay for Anonymous-DHT Receipt Line Items: A Mechanical Checklist for Digest-Bound Claims.
\newblock Synthesis~20, The-One-True-Archive (2026).


\bibitem{series:receiptacct}
Anonymous.
\newblock Receipt Accounting Rulebook: Checkpoints, Composition, and Ledger-Friendly Summaries.
\newblock Series synthesis note (Synthesis~19), 2026. (This archive.)

\end{thebibliography}

\end{document}
